\section{Análise Wireshark}

\preencher{Registrar a captura UDP durante uma transferência do arquivo
\textit{large}. Usar o filtro \texttt{udp and host 137.131.178.229}. Informar
data, duração, versão do Wireshark e arquivo de captura.}

\figura{wireshark_sessao_completa.png}{0.95\linewidth}
  {Visão geral da sessão completa: HELLO, REQs, FIN e DONE.}
  {fig:wireshark-sessao}

\preencher{Anotar a captura e explicar os principais intervalos da sessão.}

\figura{wireshark_retransmissoes.png}{0.95\linewidth}
  {Zoom em pelo menos dois eventos de retransmissão, com REQ original,
   REQ retransmitido e delta de tempo.}
  {fig:wireshark-retransmissoes}

\preencher{Identificar os dois eventos, seus números de sequência, timestamps
e deltas medidos diretamente na captura.}

\figura{wireshark_csv.png}{0.95\linewidth}
  {Correlação entre um evento do Wireshark e a linha correspondente no CSV.}
  {fig:wireshark-csv}

\preencher{Informar o timestamp do pacote, o segmento correspondente e a linha
ou registro equivalente no CSV.}

\subsection{Medição de RTT}
\begin{table}[H]
  \centering
  \caption{RTT medido no Wireshark versus RTT registrado pelo cliente.}
  \label{tab:rtt-wireshark}
  \begin{tabular}{lrrr}
    \toprule
    Segmento & RTT Wireshark (ms) & RTT CSV (ms) & Diferença (ms)\\
    \midrule
    \preencher{seq 1} & \preencher{ms} & \preencher{ms} & \preencher{ms}\\
    \preencher{seq 2} & \preencher{ms} & \preencher{ms} & \preencher{ms}\\
    \bottomrule
  \end{tabular}
\end{table}

\preencher{Comparar as medições e explicar diferenças causadas por timestamp,
processamento local, retransmissões ou pela definição de início/fim usada
pelo cliente.}
